% !PRINTED 200
\weberpage{200}
Diese Substitution muss in $G$ vorkommen und nach dem Satze 1. muss $c_1 = c_2$, also $\Theta_1$ mit $\Theta_2$ identisch sein. Die Zahlen $1, \varepsilon_1, \varepsilon_2, \ldots, \varepsilon_{\nu-1}$ müssen also zusammen alle $\nu^{\text{ten}}$ Einheitswurzeln enthalten, und also mit den Potenzen einer primitiven unter ihnen übereinstimmen. Da andererseits alle Potenzen von einer der Grössen $\Theta$ in der Gruppe (1) vorkommen müssen, so wird, wenn das in $\Theta$ vorkommende $\varepsilon$ eine primitive $\nu^{\text{te}}$ Einheitswurzel ist, die ganze Gruppe (1) durch die Potenzen von $\Theta$ erschöpft. % [sic/reference: both occurrences of ``Gruppe (1)'' are printed; the stabilizer group Q in (2), named in the next sentence, is required by the argument.]
Also besteht diese Gruppe $Q$ aus den Elementen

\weq{4}{1, \Theta, \Theta^2, \ldots, \Theta^{\nu-1}.}

Nehmen wir (nach dem Satze 2.) $\Theta$ als multiplicative Substitution an, so ist 0 der zweite Pol von $\Theta$ und zugleich der zweite Pol der sämmtlichen Substitutionen (4). Es ist daher dieser Pol gleichfalls ein mindestens \gesperrt{$\nu$-zähliger}. Da aber beide Pole von $\Theta$ in dieser Betrachtung vertauscht werden können, so erhalten wir die Sätze:

3. \gesperrt{Ein $\nu$-zähliger Pol bestimmt eine in $G$ enthaltene cyklische Gruppe $Q$ vom $\nu^{\text{ten}}$ Grade.}

4. \gesperrt{Die beiden Pole irgend einer Substitution $\Theta$ der Gruppe $G$ sind gleichzählige Pole und sind die beiden einzigen Pole der Gruppe $Q$.}

Nach \S.\,2 lassen sich $\mu - 1$ Elemente

\weq{5}{\psi_1, \psi_2, \ldots, \psi_{\mu-1}}

in $G$ so auswählen, dass $\psi_1 Q, \psi_2 Q, \ldots, \psi_{\mu-1} Q$ die Nebengruppen zu $Q$ sind, und dass also

\weq{6}{G = Q + \psi_1 Q + \psi_2 Q + \cdots + \psi_{\mu-1} Q}

wird. Wir setzen nun

\weq{7}{a_1 = \psi_1(a), a_2 = \psi_2(a), \ldots, a_{\mu-1} = \psi_{\mu-1}(a).}

Die Grössen $a_1, a_2, \ldots, a_{\mu-1}$ sind nicht nur alle von $a$, sondern auch unter einander verschieden. Denn wäre etwa

\[ \psi_1(a) = \psi_2(a), \]

so würde folgen:

\[ a = \psi_1^{-1}\psi_1(a) = \psi_1^{-1}\psi_2(a), \]

und $\psi_1^{-1}\psi_2$ wäre in $Q$ enthalten, also $\psi_2$ in $\psi_1 Q$, was gegen die Voraussetzung ist. Es ist nun leicht nachzuweisen, dass diese Werthe $a_1, a_2, \ldots, a_{\mu-1}$ sämmtlich $\nu$-zählige Pole der Gruppe sind. Es genügt, dies für $a_1$ zu zeigen.
